\documentclass[11pt,a4paper]{article}
\usepackage[utf8]{inputenc}
\usepackage[spanish]{babel}
\usepackage{amsmath, amsfonts, amssymb}
\usepackage{graphicx}
\usepackage{hyperref}
\usepackage{xcolor}
\usepackage{geometry}
\usepackage{dirtree}
\usepackage{tikz}
\usetikzlibrary{shapes.geometric, arrows.meta, positioning}
\usepackage{booktabs}
\usepackage{fancyhdr}
\usepackage{listings}

% Configuración de márgenes
\geometry{left=2.5cm,right=2.5cm,top=3cm,bottom=3cm}

% Configuración de Fancyhdr (Encabezado y pie de página)
\pagestyle{fancy}
\fancyhf{}
\fancyhead[L]{\textbf{Pipeline Filogenético SARS-CoV-2}}
\fancyhead[R]{\nouppercase{\leftmark}}
\fancyfoot[C]{\thepage}
\renewcommand{\headrulewidth}{0.4pt}
\renewcommand{\footrulewidth}{0.4pt}

% Configuración de colores para lstlisting
\definecolor{codegray}{rgb}{0.5,0.5,0.5}
\definecolor{codepurple}{rgb}{0.58,0,0.82}
\definecolor{backcolour}{rgb}{0.96,0.96,0.94}

\lstdefinestyle{mystyle}{
    backgroundcolor=\color{backcolour},   
    commentstyle=\color{codegray}\itshape,
    keywordstyle=\color{blue}\bfseries,
    numberstyle=\tiny\color{codegray},
    stringstyle=\color{codepurple},
    basicstyle=\ttfamily\footnotesize,
    breakatwhitespace=false,         
    breaklines=true,                 
    captionpos=b,                    
    keepspaces=true,                 
    numbers=left,                    
    numbersep=10pt,                  
    showspaces=false,                
    showstringspaces=false,
    showtabs=false,                  
    tabsize=2,
    frame=single,
    rulecolor=\color{black!30}
}
\lstset{style=mystyle}

% Estilos para el diagrama TikZ
\tikzset{
    process/.style={rectangle, rounded corners, minimum width=4cm, minimum height=1cm, text centered, draw=black, fill=blue!10, font=\sffamily\small\bfseries},
    arrow/.style={thick,->,>=stealth}
}

% Título y autor
\title{\textbf{Arquitectura para Inferencia Filogenética de SARS-CoV-2}}
\author{Guillermo Gutiérrez Caracuel}
\date{\today}

\begin{document}

\maketitle
\thispagestyle{empty}

\begin{abstract}
    Este documento describe un pipeline automatizado de alto rendimiento para la inferencia filogenética de secuencias de SARS-CoV-2. Orquestado mediante \textit{Snakemake}, el flujo de trabajo integra herramientas de bajo nivel en Rust para filtrado inicial, alineamiento paralelo con MAFFT, control de calidad matricial en Python (NumPy/Biopython) y construcción de árboles filogenéticos por Máxima Verosimilitud (\textit{Maximum Likelihood}) usando IQ-TREE.
\end{abstract}

\vspace{0.5cm}

\section{Introducción}
El análisis filogenético a gran escala de secuencias virales requiere soluciones eficientes y completamente automatizadas. El objetivo principal de este proyecto es proporcionar una \textbf{arquitectura de automatización y alto rendimiento} capaz de procesar de manera masiva secuencias genómicas del virus SARS-CoV-2. Esto permite generar árboles filogenéticos robustos con una intervención manual mínima, garantizando la reproducibilidad técnica, la escalabilidad estructural y una gran eficiencia computacional.

\section{Arquitectura del Sistema}
Para asegurar la mantenibilidad y modularidad del software, el repositorio sigue una jerarquía de directorios bien definida que separa los datos crudos, el código fuente y los artefactos generados.

\vspace{0.5cm}
\dirtree{%
    .1 sars-cov-2-phylogeny/.
    .2 bin/ \dots{} Binarios precompilados (\texttt{iqtree3}, \texttt{sars\_filter}, \texttt{datasets}).
    .2 data/ \dots{} Almacenamiento de secuencias FASTA crudas y procesadas.
    .2 doc/ \dots{} Documentación académica y técnica.
    .2 results/ \dots{} Resultados finales (alineamientos, árboles y gráficos QC).
    .2 sars\_filter/ \dots{} Código fuente en Rust para el filtrado rápido.
    .2 scripts/ \dots{} Scripts de análisis vectorial en Python.
    .2 Snakefile \dots{} Grafo de ejecución del orquestador.
}
\vspace{0.5cm}

\section{Diseño del Grafo Acíclico Dirigido (DAG)}
El flujo de ejecución está gobernado por \textbf{Snakemake} a través de un Grafo Acíclico Dirigido (DAG), el cual asegura que cada regla se ejecute únicamente cuando sus dependencias directas e indirectas estén satisfechas.

\begin{figure}[htbp]
    \centering
    \begin{tikzpicture}[node distance=2cm]
        \node (descarga) [process] {1. Descarga (NCBI)};
        \node (rust) [process, below of=descarga] {2. Rust Filter (\texttt{sars\_filter})};
        \node (mafft) [process, below of=rust] {3. MAFFT Alignment};
        \node (python) [process, below of=mafft] {4. Python QC (NumPy)};
        \node (iqtree) [process, below of=python] {5. IQ-TREE Inference};

        \draw [arrow] (descarga) -- (rust);
        \draw [arrow] (rust) -- (mafft);
        \draw [arrow] (mafft) -- (python);
        \draw [arrow] (python) -- (iqtree);
    \end{tikzpicture}
    \caption{Grafo Acíclico Dirigido (DAG) que ilustra las 5 fases secuenciales del pipeline.}
    \label{fig:dag_pipeline}
\end{figure}

\section{Fundamentos Matemáticos e Implementación}

\subsection{Filtrado de Secuencias (Rust)}
Para lidiar con el alto volumen de datos genómicos, el primer paso de limpieza se realiza empleando el binario \texttt{sars\_filter} implementado en Rust. Este algoritmo nativo examina iterativamente cada secuencia del genoma viral con una \textbf{complejidad algorítmica de $\mathcal{O}(n)$}, siendo $n$ el número de secuencias analizadas.

El filtro computacional asegura que la secuencia posea una longitud biológicamente plausible en un genoma libre de deleciones críticas, definida estrictamente en el intervalo:
\begin{equation}
    L \in [29000, 30500] \text{ pares de bases}
\end{equation}

\subsection{Inferencia Filogenética (IQ-TREE)}
La fase final reconstruye la historia evolutiva de las secuencias utilizando la metodología de \textbf{Máxima Verosimilitud} (\textit{Maximum Likelihood}). Bajo este riguroso marco estadístico, IQ-TREE busca la topología del árbol $T$ y los parámetros del modelo evolutivo $\theta$ que maximicen la probabilidad de observar el alineamiento de los datos reales $D$. La función de verosimilitud se define matemáticamente como:
\begin{equation}
    \mathcal{L}(\theta, T) = P(D \mid \theta, T)
\end{equation}
Para determinar de forma no sesgada el modelo evolutivo óptimo $\theta$, se emplea internamente la herramienta ModelFinder, la cual evalúa distintos modelos de sustitución nucleotídica (ej. GTR+F+I) penalizando la sobreparametrización mediante el \textit{Criterio de Información de Akaike} (AIC) o el \textit{Criterio de Información Bayesiano} (BIC).

\section{Reproducibilidad y Entorno de Ejecución}
La portabilidad técnica y la consistencia del pipeline a través de diferentes máquinas (basadas en Linux x86\_64) se logra mediante herramientas de aislamiento de software como \texttt{uv} y el suministro de archivos binarios estáticos. A continuación, se detallan los comandos necesarios para reproducir el entorno:

\begin{lstlisting}[language=bash, caption={Instrucciones de orquestación y arranque del pipeline}]
# 1. Sincronizar el entorno estricto de dependencias de Python
uv sync

# 2. Otorgar permisos de ejecucion a los binarios locales
chmod +x bin/sars_filter bin/iqtree3 bin/datasets

# 3. Lanzar la topologia DAG de Snakemake (ej: empleando 16 hilos)
uv run snakemake -c 16
\end{lstlisting}

Una vez finalizado el subproceso automatizado, los ficheros resultantes (\texttt{.treefile} en formato Newick y los gráficos de calidad poblacional en \texttt{.pdf}) quedarán localizados en el directorio \texttt{results/}.

\end{document}
